\documentclass[12pt, titlepage]{article}

\usepackage{amsmath, mathtools}

\usepackage[round]{natbib}
\usepackage{amsfonts}
\usepackage{amssymb}
\usepackage{graphicx}
\usepackage{colortbl}
\usepackage{xr}
\usepackage{hyperref}
\usepackage{longtable}
\usepackage{xfrac}
\usepackage{tabularx}
\usepackage{float}

\usepackage{booktabs}
\usepackage{multirow}
\usepackage[section]{placeins}
\usepackage{caption}
\usepackage{fullpage}
\usepackage[table]{xcolor}
\usepackage{array} 
\usepackage{colortbl}
\usepackage{makecell}
\usepackage{longtable}



\definecolor{linkblue}{RGB}{0, 102, 204}
\hypersetup{
    colorlinks=true,
    linkcolor=linkblue,
    urlcolor=linkblue,
    citecolor=linkblue,
    filecolor=linkblue
}


\title{Usability Report}
\author{Maham S}
\date{\today}

\begin{document}
\maketitle
\newpage
\tableofcontents
\newpage

\section{Introduction}
Evaluation of qualitative aspects of the Physio Companion app was done through a usability study.
The project timeline and resources limit the participant pool size, and as a result,
this study does not provide any quantitave insights.
This document outlines the testing methodology and concludes with the resulting findings. \\
The goals of the study are: 
\begin{itemize}
    \item Assess the ease of use of the app across a diverse demographic
    \item User understanding of the terminology within the app
    \item User perception of the flow and navigation within the app
    \item Identification of any friction points 
    \item Feedback on additional and beneficial features within the app
\end{itemize}

\section{Techniques}
The techniques used to conduct these tests were a mixture of both in-person and remote moderated testing. Data was
collected via a questionnaire or an interview. The choice was up to the discretion of the moderator, but both formats covered
the same core topics.
They were first asked a series of questions inquiring some background information. 
The participants were then shown how to install the app on their smartphone. After successful installation of the program, the facilitator prompted
the participant on the app's use. An additional investigator was playing the role of an observer recording the user's behaviour. 
At the end of the study, participants were asked for their final impressions. The study was voluntary and lasted approximately 30-40 minutes.

\section{Methodology}
In order to gain a deeper understanding of the user experience of the application, a series of questionnaires were constructed
for each stage of testing (\href{Pre-sessionQuestionnaire.pdf}{pre-session}, \href{run:SessionQuestionnaire.pdf}{session}, 
and \href{Post-sessionQ.pdf}{post-session}). In addition, the 'session' questionnaire was further subdivided
in order to examine the different main actions of the application. This breakdown is crucial
to understanding the different facets of the application. Obtaining feedback promptly after completing each section will 
allow immediate and reliable reporting of the participants' impressions. The questionnaires were constructed as a guide and 
investigators were free to use either an interview-style approach or use the explicit questionnaire. If the former was used, 
it was essential that the topics mentioned in the questionnaire were addressed. 



\section{Participant profiles}
The participant pool comprises of individuals of varying ages, genders and occupations, which ensures coverage of 
a broad target user base. The coverage details are outlined below:
\begin{itemize}
    \item Name
    \item Age: Participants of various ages were chosen, with the sole restriction being that they had to be over 18.
    \item Gender: In an ideal case, an equal and all-encompassing balance of genders would be sought out. This would include 2 female 
    or male identifying individuals and one non-binary individual.
    \item Person with disability: Another criteria that ensures diversity is enlisting participants with disabilities. This will aid in
    identifying any accessibility barriers that may go unnoticed. In an ideal case, at least 1 or 2 participants would be those who have a disability.
    \item Occupation: This criteria was selected to give insight into lifestyle factors. For the purposes of this study, any bias towards conducting 
    it for the same occupation would skew the results.
    \item Level of digital literacy: This measure pertains to smartphone usage. It is divided into low, medium, or high expertise indicating 
    how comfortable the user is operating the device.
    \item Familiarity with injury to the knee: This relates to how familiar the participant is with a knee injury, whether they 
    are familiar with it through their occupation (e.g. orthopaedic surgeon, physiotherapist) or whether they have had a knee injury in the past.
\end{itemize}

\textit{Participant 1}
\begin{itemize}
    \item Age: Late 20s
    \item Gender: Male
    \item Person with disability: No
    \item Occupation: Physiotherapist
    \item Level of digital literacy: High
    \item Familiarity with injury to the knee: Very familiar
\end{itemize}

\textit{Participant 2}
\begin{itemize}
    \item Age: Late 20s
    \item Gender: Male
    \item Person with disability: No 
    \item Occupation: Physiotherapist
    \item Level of digital literacy: High
    \item Familiarity with injury to the knee: Very familiar
\end{itemize}

\textit{Participant 3}
\begin{itemize}
    \item Age: early 60s
    \item Gender: Male
    \item Person with disability: No
    \item Occupation: Technician
    \item Level of digital literacy: High
    \item Familiarity with injury to the knee: Familiar
\end{itemize}

\textit{Participant 4}
\begin{itemize}
    \item Name
    \item Age
    \item Gender
    \item Person with disability
    \item Occupation
    \item Level of digital literacy
    \item Familiarity with injury to the knee
\end{itemize}

\textit{Participant 5}
\begin{itemize}
    \item Age: Mid 20s
    \item Gender: Female
    \item Person with disability: No
    \item Occupation: Geoscientist
    \item Level of digital literacy: High
    \item Familiarity with injury to the knee: Familiar
\end{itemize}
Due to project timeline constraints, the participant pool was limited in size, and posed challenges in recruiting 
individuals who satisfied the criteria mentioned above. Furthermore, it was not possible to recruit a participant 
with a knee injury. However, there were some that had undergone physiotherapy treatment for other health reasons. \\
Two of the participants were physiotherapists. They were recruited to assess the PT flow of the application and 
provide their perspective.The remaining participants varied in age and fulfilled different combinations 
of the criteria above. 

\section{Testing}
All questionnaires' accompanying logic is explained \href{QuestionsExplanations.pdf}{here}.
\subsection{Pre-session questionnaire}
See \href{Pre-sessionQuestionnaire.pdf}{pre-session questionnaire}.
\subsection{Session questionnaire}
See \href{SessionQuestionnaire.pdf}{session questionnaire}. 
\subsection{Post-session questionnaire}
See \href{Post-sessionQ.pdf}{post-session questionnaire}.


\section{Results}
 %https://aging.jmir.org/2023/1/e43197/ 
 % Use this to talk about aging and tech 
%There were 5 participants of the study: 2 PTs and 3 others that fulfilled the criteria mentioned above. \\

\noindent \textbf{Physiotherapists} \\
These points emerged as common feedback for both of the physiotherapist participants (See Appendix A). 
\begin{itemize}
    \item Both participants emphasized the inclusion of an RPE (Rate of Perceived Exertion) from the patient
    as an additional measure.
    \item They also suggested incorporated the number of sets performed in each exercise session on the 
    Feedback tab.
\end{itemize}

\section{Discussion}
Usability testing proved to be an invaluable asset in evaluating the application. The results were 
enlightening and were immensely beneficial to the development team in refining the final product. \\
For the physiotherapist participants, the study was conducted in a conversational interview format, because it proved to be 
more natural than a questionnaire. 
The physiotherapists 
tested the PT flow of 
the application and highlighted some 
improvements (see above). 
They noted that the patient's subjective experience was an indespensable measurement, which is particularly 
valuable when determining exercise modifications. They provided examples of how this may indicate how difficult
or easy the exercise is for the patient and whether their performance is impacted by other limitations (eg. pain).
Including the number of sets is useful for the same reason. Both participants found the navigation and
its use within the application to be straightforward. Their positive response to the application
suggests that it has strong potential and can be advantageous to patient recovery. 


\section{Conclusion}

\section{Appendix A}
\subsection{Usability Test Feedback: Physiotherapist 1}

\begin{itemize}
    \item \textbf{On a scale of 1 to 10, how easy was it for you to log in?} \\
    Response: "Pretty good. Like, no complaints. Yeah, pretty straightforward."

    \item \textbf{Tell me what you see and what you think is going on in this page [the homepage]?} \\
    Response: Recognized it as a welcome page. Correctly identified that linked patients are listed on the left and the patient activity history shows recent exercises.

    \item \textbf{From this page, can you tell me what an "invite code" is?} \\
    Response: "No." The moderator then explained it is for linking patients to the physiotherapist's organization.

    \item \textbf{Was there any sort of confusion when you initially landed on the homepage, and what would make it easier to understand?} \\
    Response: The "patient activity history" on the main dashboard feels slightly repetitive since that information is also available in the feedback section. He would prefer to swap out the activity history on the main dashboard with a direct list of the specific exercises assigned to that patient, as that is what physiotherapists look at the most.

    \item \textbf{Is there a specific organizational structure that would help you to better address your patients (e.g., categorizing by most recent patient, submission, or exercise)?} \\
    Response: The current format looks good, but he recommended removing the "Recommended" versus "Optional" categorization. Exercises should be treated like prescribed medication—patients just have to do them. It is best to keep it simple and just list the required exercises.

    \item \textbf{Do you feel like you have enough information in the exercise logs to make decisions or give comments about a patient's progress?} \\
    Response: Requested the addition of "Sets" and "Reps" (e.g., 10 reps for 2 sets), as well as "Holds" (hold duration for a stretch). This provides a much more realistic picture of what the patient is doing.

    \item \textbf{Is there anything particular that you found was hard to do when modifying the instructions?} \\
    Response: No, but having a physically separate input box for assigning "Reps" and "Sets" (rather than just typing it into a general instruction text box) would be perfect.

    \item \textbf{If we added a second exercise, would you prefer a list of every available exercise or an "add exercise" dropdown?} \\
    Response: An "add exercise" dropdown, if possible.

    \item \textbf{How useful would you say are the charts and graphs for tracking a patient's progress over time?} \\
    Response: Mildly helpful. He noted that he currently goes mostly by what the patient tells him. However, he strongly recommended keeping it in the app because it's a great extra feature that helps the app stand out from competitors.

    \item \textbf{Was there any data that you expected to see here that is missing?} \\
    Response: It would be highly helpful to track patient effort levels, specifically RPE (Rate of Perceived Exertion), and pain levels on a scale of 0 to 10. This is especially crucial for post-op patients to ensure they aren't pushing too hard and causing excruciating pain.

    \item \textbf{Do you feel that it's better to remove all patient history once a user is deleted, or have a log of previous patients?} \\
    Response: By law, healthcare professionals must keep medical records for about 8 years. Physiotherapists generally never delete patients.

    \item \textbf{Follow-up Question: Would it be safe to say we can have an "offload" or "inactive" feature instead of "delete"?} \\
    Response: Yes. Labeling it as an "inactive account" or "discharged" is a good idea.

    \item \textbf{On a scale of 1 to 10, how easy was the app to navigate without any instructions?} \\
    Response: A solid 8.5.

    \item \textbf{What was the most confusing part of the app?} \\
    Response: The feedback section wasn't super obvious at first glance. He suggested adding a quick blurb or context at the top of the screen, such as: "Click on an exercise to leave your patients a comment."

    \item \textbf{If there was a feature you would add or change to make this more useful in your practice, what would that be?} \\
    Response: An easy way to add and remove exercises. Physiotherapists spend 99\% of their time prescribing and modifying exercises, and it is the biggest hassle at the end of the day.

    \item \textbf{For tracking, is it required for the app's camera to automatically measure sets/holds, or is tracking total completed reps alongside the requested reps/sets enough?} \\
    Response: Displaying the requested reps/sets and only having the app count the completed reps is perfectly fine. The app does not need to automatically measure hold lengths or sets.
\end{itemize}
\subsection{Usability Test Feedback: Physiotherapist 2}

\begin{itemize}
    \item \textbf{On a scale of 1 to 10, how easy was it for you to log in?} \\
    Response: Pretty straightforward. 

    \item \textbf{Tell me what you see and what you think is going on in this page [the homepage]?} \\
    Response: Explained the main dashboard correctly. Noted that the patients listed are his patients, and the history section is to see whether they are doing their assigned physio. He assumed that the invite code was to be sent to a new patient to add them to this application.

    \item \textbf{From this page, can you tell me what an "invite code" is?} \\
    Response: Assumed correctly—the invite code is to be given to a new patient, who can then enter it on their side to join the physiotherapist's planning.

    \item \textbf{Was there any sort of confusion when you initially landed on the homepage, and what would make it easier to understand?} \\
    Response: Overall, it was pretty straightforward. He did note, however, that clicking a patient to view their history was not immediately obvious—he did not do that himself without prompting.

    \item \textbf{Is there a specific organizational structure that would help you to better address your patients (e.g., categorizing by most recent patient, submission, or exercise)?} \\
    Response: Looks pretty good. I think adding more information would be helpful (e.g. frequency per day, sets)

    \item \textbf{Is there anything particular that you found was hard to do when modifying the instructions?} \\
    Response: No, it was pretty self-explanatory. He suggested that adding pictures or a visual component would make it easier on their end. Additionally, he reiterated that including sets, reps, and holds would be beneficial.

    \item \textbf{If we added a second exercise, would you prefer a list of every available exercise or an "add exercise" dropdown?} \\
    Response: "I think having some sort of search bar would be easier to navigate through the exercises."

    \item \textbf{How useful would you say are the charts and graphs for tracking a patient's progress over time?} \\
    Response: He emphasized the need to track not just sets, but also which specific exercise was performed. For example, if the rep count is 30, it is unclear whether that was a 3x10 or a 2x15, or what exercise it corresponds to. He also suggested including RPE (Rate of Perceived Exertion). \\
    He placed a big emphasis on distinguishing between multiple exercises (e.g., upper body vs. lower body), noting that an overall rep count becomes less meaningful when exercises are mixed.

    \item \textbf{Was there any data that you expected to see here that is missing?} \\
    Response: He highlighted the value of granular data—for example, if a patient did 3x10 and the first two sets went well but on the third set, the sixth rep was difficult, flagging that specific moment would be more important to a physio. This level of detail allows them to tailor the exercise regimen more effectively (e.g., identifying if the load is too much or too little).

    \item \textbf{Do you feel like you have enough information in the exercise logs to make decisions or give comments about a patient's progress?} \\
    Response: "Not really." While objective numbers are nice, having subjective comments from the user would be much more helpful from a physio's point of view. Again, he mentioned that tracking RPE and how the patient felt during exercises would be valuable.

    \item \textbf{Do you feel that it's better to remove all patient history once a user is deleted, or have a log of previous patients?} \\
    Response: Deleting is fine, so long as there is some sort of export function. He expressed his frustration with how long charting currently takes. It would be ideal to export patient notes (e.g., as a PDF) directly into their record-keeping application. This way, the practice could keep records without needing to retain the patient's account indefinitely.

    \item \textbf{On a scale of 1 to 10, how easy was the app to navigate without any instructions?} \\
    Response: Overall, it was pretty straightforward, and the UI was easy to use. However, he noted that within the feedback section, adding a visual indicator (e.g., a "\textgreater" icon) would help show that tiles are clickable—he had to be told they were interactive.

    \item \textbf{What was the most confusing part of the app?} \\
    Response: Mentioned above (clickable tiles were not immediately obvious).

    \item \textbf{If there was a feature you would add or change to make this more useful in your practice, what would that be?} \\
    Response: Mentioned above (export functionality, visual indicators for clickable elements, and more granular exercise tracking).
\end{itemize}

\end{document}